%Preamble
\documentclass[12pt,a4paper]{article}

\usepackage[T1]{fontenc}
\usepackage[utf8]{inputenc}
\usepackage[romanian]{babel}
\usepackage{lmodern}
\usepackage[margin=2.5cm]{geometry}
\usepackage{microtype}

\setlength{\parindent}{0pt}
\setlength{\parskip}{0.8em}


\title{\textbf{Titlu: Agenți inteligenți pentru medii virtuale} }

\begin{document}
\maketitle

\textbf{Student:} Ing.\ Marius-Vlad DUMITRU

\section{\textbf{Descrierea temei} }

    Proiectul explorează dezvoltarea și integrarea agenților inteligenți autonomi în medii virtuale cu accent pe: percepție, luarea deciziilor, acțiune în conformitate cu un set de reguli predefinite.

    Acest demers contribuie la avansarea inteligenței artificiale, în special în domeniul luării deciziilor în timp real și al învățării adaptive. Prin analiza comportamentului unui agent inteligent într-un mediu virtual, proiectul își propune să abordeze limitările actuale ale modelelor de învățare profundă, cum ar fi interpretabilitatea și explicabilitatea deciziilor luate.

    Rezultatele vor fi aplicate în cadrul jocului video MicroRTS, cu scopul de a demonstra aplicabilitatea și eficiența agenților în scenarii interactive.

    Agenții vor reprezenta jucători virtuali ai jocului MicroRTS.

    Mediile virtuale for reprezenta hărți ale jocului MicroRTS, în diverse configurații, la care se adaugă diferite constrângeri impuse agenților.

    În procesul de dezvoltare se vor folosi limbajele de programare Python și Java.

    \pagebreak

\section {\textbf{Contribuția studentului} }
    \begin{enumerate}
        \item Studentul va implementa modulele necesare pentru asigurarea comunicației dintre agenții inteligenți autonomi (jucătorii) și jocul MicroRTS precum și a modulelor necesare respectării constrângerilor impuse. Arhitectura acestora va fi ajustată pentru a îmbunătăți performanța sau pentru a răspunde cerințelor specifice experimentelor.

        \item Studentul va implementa modificări asupra codului sursă al jocului MicroRTS, acolo unde sunt necesare pentru asigurarea interacțiunii dintre joc și agenți dar și pentru facilitarea cerințelor specifice experimentelor.

        \item Studentul va antrena modele pentru agenții autonomi folosind seturi de date relevante, iar în cazul în care acestea nu sunt suficiente, va explora metode pentru generarea sau extinderea lor.

        \item Studentul va efectua experimente prin modificarea unor componente de arhitectură sau practici de antrenare pentru a evalua impactul acestora asupra performanței agenților, incluzând timpii de decizie și rata de succes în sarcini complexe.

        \item Validarea rezultatelor obținute va include o comparație cu alte metode existente din literatura de specialitate, iar eficiența sistemului propus va fi evaluată prin metrici specifice pentru performanța agenților în medii virtuale.
    \end{enumerate}

\begin{thebibliography}{9}
    \setlength{\itemsep}{0.5em}

    \bibitem{liu2024} R.-Z. Liu, Y. Shen, Y. Yu and T. Lu, ``Revisiting of AlphaStar,'' in \textit{IEEE Transactions on Games}, vol.~16, no.~2, pp.~317--330, June 2024, doi: 10.1109/TG.2023.3265975.

    \bibitem{tao2023} J. Tao et al., ``Explainable AI for Cheating Detection and Churn Prediction in Online Games,'' in \textit{IEEE Transactions on Games}, vol.~15, no.~2, pp.~242--251, June 2023, doi: 10.1109/TG.2022.3173399.

    \bibitem{singh2024} J. Singh, S. Rani and G. Srilakshmi, ``Towards Explainable AI: Interpretable Models for Complex Decision-making,'' 2024 International Conference on Knowledge Engineering and Communication Systems (ICKECS), Chikkaballapur, India, 2024, pp.~1--5, doi: 10.1109/ICKECS61492.2024.10616500.

    \bibitem{li2024} W. Li, W. Liu, S. Shao, S. Huang and A. Song, ``Attention-Based Intrinsic Reward Mixing Network for Credit Assignment in Multiagent Reinforcement Learning,'' in \textit{IEEE Transactions on Games}, vol.~16, no.~2, pp.~270--281, June 2024, doi: 10.1109/TG.2023.3263013.

    \bibitem{goodfriend2024} S. Goodfriend, ``A Competition Winning Deep Reinforcement Learning Agent in microRTS,'' in \emph{2024 IEEE Conference on Games (CoG)}, Milan, Italy, 2024, pp. 1--8, doi: 10.1109/CoG60054.2024.10645610.
\end{thebibliography}

\end{document}
